\section{Language Interfaces}

External applications access Siderust through a C ABI and Python bindings. The
\texttt{siderust-ffi} crate exports C-compatible coordinate values, tagged target descriptions,
opaque resource handles and status codes. Runtime ephemeris and transformation contexts are created
and released through explicit functions. Coordinate metadata identify the frame, center and unit;
errors distinguish invalid arguments, unsupported representations, missing data and out-of-range
queries. The C header is generated from the Rust definitions with \texttt{cbindgen}, and the C++
interface is built on this ABI.

The Python implementation uses PyO3 to expose observers, directions, positions, ephemerides and
observing-event operations. Numerical routines remain in Rust. Python positions carry frame, center
and unit attributes that are checked by the bindings. The time facade reuses Tempoch's Python
package, although observation functions that take numerical MJD or JD arguments still rely on their
documented scale conventions. Coverage and representations vary by entry point; the bindings do not
expose Rust's generic parameters as host-language compile-time constraints.

The experiments in Sections~7--8 evaluate the native Rust interfaces. They provide the numerical
reference cases for future tests of the corresponding foreign-language calls; cross-language
residuals are not reported here.
